\chapter{ÜST SEVİYE TASARIM}
\thispagestyle{empty}

Sistemin üst seviye tasarımı, birbirini besleyen dört ana katmandan oluşmaktadır: Veri Hazırlama, Metin Kodlama, Graf Öğrenimi ve Öneri Sunumu.

% TODO: buraya diyagram
% \begin{figure}[H]
%     \centering
%     \includegraphics[width=\textwidth]{architecture.png}
%     \caption{Sistemin üst seviye mimarisi.}
%     \label{fig:architecture}
% \end{figure}

\section{Katman 1 -- Veri Hazırlama}

Bu katman, ham Wikipedia verisini model eğitimine uygun bir yapıya dönüştürmekten sorumludur.

\begin{enumerate}
    \item Wiki-CS veri setinin graf topolojisi (komşuluk matrisi) alınır.
    \item Düğüm başlıkları normalize edilerek Wikipedia'daki gerçek makale başlıklarıyla eşleştirilir.
    \item Makale metinleri, öncelikle Wikipedia STEM Corpus üzerinden, bulunamayanlar ise Wikipedia API aracılığıyla elde edilir.
    \item Elde edilen metinler temizlenerek bir sonraki katmana aktarılır.
\end{enumerate}

Çıktı: 11.701 düğümlü yönlü graf (komşuluk matrisi) ve her düğüme karşılık gelen makale metni.

\section{Katman 2 -- Metin Kodlama}

Bu katmanda, her makale metni ön eğitimli cümle gömme modelleri aracılığıyla yoğun vektörlere dönüştürülür. İki model karşılaştırmalı olarak değerlendirilmektedir:

\begin{itemize}
    \item BAAI/bge-m3 $\rightarrow$ 1024 boyutlu gömme vektörü
    \item Alibaba-NLP/gte-modernbert-base $\rightarrow$ 768 boyutlu gömme vektörü
\end{itemize}

Bu vektörler, graf üzerindeki düğüm öznitelikleri (node features) olarak kullanılacaktır.

Çıktı: Düğüm öznitelik matrisi $\mathbf{X} \in \mathbb{R}^{N \times d}$ (burada $N = 11.701$, $d \in \{768, 1024\}$).

\section{Katman 3 -- Graf Öğrenimi (GAT)}

Bu katman, metin gömme vektörlerini graf yapısıyla birleştirerek zenginleştirilmiş düğüm temsillerini öğrenir.

\begin{enumerate}
    \item Düğüm öznitelik matrisi $\mathbf{X}$ ve komşuluk matrisi $\mathbf{A}$, GAT modeline girdi olarak verilir.
    \item GAT'ın çok başlıklı dikkat mekanizması (multi-head attention), her kenar için önem ağırlığını hesaplar.
    \item Birden fazla GAT katmanı ardışık olarak uygulanarak, her düğüm çok adımlı komşuluk bilgisini içeren zenginleştirilmiş bir temsil elde eder.
    \item GCN aynı girdi üzerinde baseline model olarak ayrıca eğitilir ve sonuçlar karşılaştırılır.
\end{enumerate}

Çıktı: Zenginleştirilmiş düğüm gömme matrisi $\mathbf{Z} \in \mathbb{R}^{N \times d'}$, burada her vektör hem makalenin anlamsal içeriğini hem de graf yapısındaki konumunu yansıtır.

\section{Katman 4 -- Öneri ve Hiyerarşik Sunum}

Bu katman, kullanıcı sorgusunu alarak nihai öneriyi üretir.

\begin{enumerate}
    \item Kullanıcı sorgusu, aynı cümle gömme modeli ile vektörleştirilir.
    \item Sorgu vektörü ile düğüm gömme vektörleri ($\mathbf{Z}$) arasında benzerlik hesaplanarak (kosinüs benzerliği) en yakın $k$ makale aday olarak belirlenir.
    \item Aday makaleler, yerel olarak çalıştırılan bir LLM'e (DeepSeek-R1 veya benzeri) iletilir.
    \item LLM, aday makaleleri kavramsal ön koşul ilişkilerine göre hiyerarşik bir öğrenme patikası olarak düzenler ve kullanıcıya sunar.
\end{enumerate}

Çıktı: Kullanıcıya sunulan, ``nereden başlamalıyım'' sorusuna yanıt veren yapılandırılmış bir makale öğrenme yolu.